\section{Production and utility in an agent context \label{sec:produtility_agents}}

\begin{comment}
Contrast to what agent firms need to do; data, meaning�
Cf. Baumol: faced with a particular reality, what sort of theory is actually going to help firms succeed?
"It is at least possible that shear business acumen and experience permit management and other economic units to arrive at decisions which come close to being optimal. Moreover, in business, competition may soon eliminate firms whose decision-making is consistently poor. To the extent that these assertions are valid, optimality analysis should serve as a relatively good predictor of economic behaviour; that is, it should provide a reasonably good explanation of actual economic decisions and activities." [5]
Plus, randomness, sequence, luck may also eliminate a certain number. 
Cf. nature of optimisation; function standing in for process of groping, etc.

\end{comment}

%%%%%%%%%%%%%%%%%%%%%%%%%%
% OUTLINING UTILITY BEING USED IN MY MODELS
%%%%%%%%%%%%%%%%%%%%%%%%%%
\subsection{Using utility in an agent model \label{sec:rationality}}

In the model framework developed for this thesis (described in detail in chapter \ref{chapter:modelfw}), agents use a set of `standard' utility functions. The functions themselves have been covered in the previous chapter. 

Many agent-based critiques of economics argue that the utility-maximising approach is an unrealistic straight-jacket... blah. This section argues in support of using utility. If understood properly, and used with awareness of its limitations and meaning, utility-maximising agents are a very useful tool.

One of the central criticisms of utility theory is that people are simply not like that. (x-ref Edmonds.) This is examined in more depth below, but it is useful to quickly explain the thesis' position on this: people are not utility-maximising automatons, but utility theory \emph{is} a useful way to analyse how people react to changes in costs. There is no need to make any sophisticated case for the psychological basis of utility; it can still be used as a modelling tool. Neither does this mean more sophisticated models of behaviour are unimportant: utility is merely a simple, effective way of gaining insights into the impact of cost changes. There are some pitfalls to be avoided when interpreting utility-based agents; these are discussed below.




Some things I need to cover:

utility as tautology / hedonistic / actual internal description of psychological state. Problem of information, uncertainty. The point about realism.

What's the value of realism?

A defence of the use of utility requires understanding what the type of agent model developed here is being used for -  what it is being asked to do. [does it? That's nice.]



The basic point is that utility provides a relatively simple way of describing how people react to changes in cost. It does this job admirably - but problems arise when.. when what? 


In theory, no more need be said on the matter; in practice, more justification is required, since much confusion arises from its use.


Bouncing off Einstein's almost cliched quote; the surprise that nature should be so describable by maths. The description has so much veracity, it's possible to forget that it is 

The tight fit, the close isomorphism, between reality and physical models means we trust those models to predict, we trust them with our lives. But it still hides the basic truth that physical models are still only a description, a way of representing reality. Map and territory are not the same thing. Forgetting that fact causes much confusion when it comes to modelling in the social sciences. In particular - and this is Friedman's point - there is a belief that physical models closely match actual structures in reality and that's where their power comes from. Of course that isn't true: they massively simplify reality, for example by treating planets as points. Adding more `realism' makes many models worse, not better. 

This is a problem for \gls{ABM} because, oddly, their ridiculousness is more visible. In reality, firms do not crawl about on a 2D plane, people do not hop effortlessly from place to place searching for the most utility. 

\subsection{Very brief history of utility}


\subsection{Is utility `credible?'}

The work of Moss and Edmonds is a good example of a general view within agent modelling: the `standard' economic approach to modelling people using utility is simply not `credible'.

1. Thomas Brenner, Computational Techniques for Modelling Learning in Economics (Springer, 1999).
That's got this PDF: http://cfpm.org/cpmrep33.html. Also in zotero. 

See also: 




\subsection{Utility as empty tautology}

Compare two examples of the use of utility from the literature: one from transport economics, the other from rational choice theory. Approaches to utility in transport economics have used the idea of revealed preference successfully to say useful things about people's choices. Button points out that - 

"the general conclusion about the idea that some overall budget mechanism governs individual travel decisions, however, must be that, to date, the evidence available still leaves many questions unanswered and the theory is still largely unproved." [92]

Yet, despite this, many directly practicable ideas have emerged through trying to understand how people react to cost changes. One that stands out is finding revealed preferences for travel time: ``if a person chooses to pay \$x to save y minutes then he/she is revealing an implicit value of time equal to at least \$(x/y) per minute.'' This has led to - 

``savings in walking and waiting times are valued at between two and three times savings in on-vehicle time - parameters that have proved to be remarkably robust over the years.'' [104]

This is striking: how people value travel time is not simply a function of the time taken. Those willing to commute by car for an hour each way would be much less likely to walk that distance. As regards this thesis, the crucial point is that no argument need be made about the underlying mechanism or the rationality of those making a choice between walking or driving to work. There is, of course, value in attempting to understand what underpins those choices, but equally there is value in accepting and applying the finding. The underlying causes may be complex, but this does nothing to lessen the importance of the result itself.
%Survey: Walters, 'value of travel time saving used in road project evaluation: a cross-country / jurisdiction comparison', Austr trans reserach forum 92.
%Paper: Hensher - 'measurement of the valuation of travel time savings', journal of transport economics and policy, 2001
%From a Leeds study, think this might be the one Varian talks about. Quarmby, 'choice of travel mode for the journey to work: some findings', journal of transport economics and policy, 67.
So people willing to commute for an hour each way would be much less likely to walk that distance.  Utility describes but does not explain this.

Compare this to rational addiction theory...

Blaug On samuelson's revealed preference theory (RPT): ``provides a perfect example of our earlier contention that the logical distinction between `assumptions' and `implications' disappears in a perfectly axiomatized theory.'' [143]



The problem is simply enough stated: if any outcome can be described as rational, then the theory is completely unfalsifiable. 

But it depends on what you're falsifying, and what you're able to falsify with the theory. Any description of gravity that we have isn't really in a position to falsify the graviton theory. It'd need some deeper level of data to uncover what's going on. 

The physical analogy can only be taken so far, of course: in reality, there are plenty of ways of showing that people are not utility maximisers. The argument for using utility must, therefore, be relatively lightly made: it is a useful tool, as long as its limitations are kept in mind.


%MM1 20pp/23	Reasoning behind bundle approach; cf. to Fowler's 'objective demand' approach. Write up.
%%%%%%%%%%%%%%%%%%%%%%%%%%
% OUTLINING UTILITY BEING USED IN MY MODELS
%%%%%%%%%%%%%%%%%%%%%%%%%%
\subsection{Rationality n ting \label{sec:rationality}}

The goal of agents in these models is the same as those in \gls{ge} - to optimise either their output or their utility. The technical challenges of achieving this in an \gls{abm} setting, however, prompts some useful questions. This section looks more closely at the nuts and bolts of getting agents to optimise.
%The process of developing the models led to a number of unexpected issues; these are a very useful way to analyse the fundamental differences (and similarities) between agent and analytic modelling approaches.

Section \ref{sec:frameworkUtility} in chapter \ref{chapter:modelfw} explains how utility is used in the model framework. This section outlines the key ideas that led to those model choices. The central issue stems from problems using `objective' utility function optimisations of the type explained in section \ref{sec:utilityintro}. In these, a utility function is constrained by a budget. This produces a set of equations giving the `objective' optimal quantity of goods, given that budget. The approach in  \gls{ge} (and many other economic models) is to assume the objective demand implies the correct level of production to meet it, and thus the correct number and scale of firms. In terms of agents, as Leijonhufvud puts it -

%http://www.forvo.com/word/leijonhufvud/ Lay-on-huh-vud.
%constrained optimisation models attribute \emph{substantive rationality} to agents. Thus, 
\begin{quote}
``in this theory utility or profit maximisation is a statement about actual performance, not just motivation... The theory does not leave room for failures to realise the relevant optima.'' (\citealt{leijonhufvud_agent-based_2006} p.1628)
\end{quote}

%In discussing `market design' using agents, Marks points out that building artificial markets has forced programmers and theorists to move away from closed-form equilibria. (\fixref{Marks chapter in handbook}) More specifically
This is not just an abstract criticism of the constrained optimisation approach: it is a problem that needs solving in agent models using utility. Fowler's approach is to make a tactical choice to allow agents their objective demand. As he puts it, ``workers select their optimal bundle of goods without reference to the actual supply of the good'' (\citealt{fowler_finding_2010} p.10). This avoids the need to solve rationing of limited goods. It also sends ``clear signals to firms about the actual level of demand for their goods at the current price'' (ibid).

As has been argued, \cmt{where argument about conscious defence of model sillinesses?} all models will have implausible-appearing assumptions, and it is the job of the modeller to consciously acknowledge what they are, and justify the analysis that includes them. The job Fowler set himself was to stick as closely to the \gls{coremodel} as possible, and keeping objective demand is a part of that. This thesis takes a different approach. 

---
%%%%%%%%%%%%%%%%%%%%
% PRODUCTION AND EOS
%%%%%%%%%%%%%%%%%%%%
\subsection{Coordinating simple production \label{sec:coord_production}}

Section x explained that, in the \gls{coremodel}, firms have a production function that gives them increasing returns from labour input. It was also noted that those increasing returns are analysed away: changes in welfare through production come about solely through changes in the equilibrium number of firms, and thus the total number of varieties consumed. 

But Leijonhufvud's point applies here too: analytic approaches skip to the end of the process of production. They describe optima, and build that description on an explanation devoid of specific actions to get there. Whether or not this approach is justified (a point discussed more below) is beside the point when it comes to building an agent model of production: \gls{abm} applies its own restrictions here, forcing the modeller to work with specific `Firm' agents, and specifying what actions they must take in order to manage production. (x-ref Fowler discussion.) As just mentioned, firms enter and exit the market also. An agent approach must decide how to address this, even if to explicitly avoid it. 

So the question is: how to build Firm agents? What are the issues? In both the case of Firms and People, this thesis has attempted to build the simplest basic economic behaviour into them (and also agrued this is justifiable). 

The most obvious starting point is to ask, what are Firms' goals? Their over-arching aim is simple enough to state: maximise net revenue, which is gross revenue minus costs. If a Firm can stay afloat doing this, they can survive. This simple aim quickly becomes complicated when considering all the options then facing a model of Firm action. Chapter \ref{chapter:modelfw} explains in detail what options are examined in this thesis, but now is a good time to explore the over-arching issues.

Section \ref{sec:fowler} discussed Fowler's eloquent encapsulation of the co-ordination issues involved: 

``the amount a firm can offer in wages is dependent on the amount it can produce and the price it receives for its goods. These quantities are affected, in turn, by consumer's wages, which depend on which firm employs them (and whether or not they have found employment at all.)'' (\citealt{fowler_taking_2007} p.277) 

This tangle of dependencies can be broken down to see how each element in an agent model might relate to each other. Firms can be modelled in ways that avoid many coordination issues, while still helping to shed light on the modelling problem. One approach is to break the cycle of dependencies completely. By providing People with a fixed wage income every turn, issues of money circulation are avoided: firms do not need to set or pay wages, and so payment can appear and disappear again in one day without the need to circulate.
\cmt{need to x-ref with numeraire stuff / ways to solve money problem / meaning of money etc. Maybe.}
%Diagram showing this? E.g. flow vs "`in >> person >> firm >> out"'




%%%%%%%%%%%%%%%%%%
% ARE BLACK BOXES TELEOLOGICAL?
%%%%%%%%%%%%%%%%%%
%\subsection{Are black box models teleological? \label{sec:blackbox_tel}}

Hayek, in his short essay, `the use of knowledge in society' (\citealt{hayek_use_1945} p.529) criticises Schumpeter for using ``objectively given quantities of commodities impinging directly upon each other, almost, it would seem, without any intervention of human minds''. Hayek quotes Schumpeter as arguing that -

\begin{quote}
``... the possibility of a rational calculation in the absence of markets for the factors of production follows for the theorist `from the elementary proposition that consumers in evaluating (`demanding') consumers' goods \emph{ipso facto} also evaluate the means of production which enter into the production of these goods.' '' (\citealt{schumpeter_capitalism_2010} p.175 in Hayek, ibid.)
\end{quote}

\cmt{this crosses over with args about monocomp, e.g. in Fowler section.}

The first thing to note is that this is an argument made by \gls{ge}: production is assumed as the mirror image of consumption, and everything else about firms, including their equilibrium number, flows from that. \footnote{There is a classic Hayek quote here: ``the problem is thus in no way solved if we can show that all the facts, if they were known to a single mind (as we hypothetically assume them to be given to the observing economist), would uniquely determine the solution; instead we must show how a solution is produced by the interactions of people each of whom possesses only partial knowledge. To assume all the knowledge to be given to a single mind in the same manner in which we assume it to be given to us as the explaining economists is to assume the problem away and to disregard everything that is important and significant in the real world..'' (\citealt{hayek_use_1945} p.530; see e.g. \citealt{miller_complex_2007} p.87)}; the important part for the argument being made is that Hayek believes Schumpeter to be ``disregard[ing] everything that is important and significant in the real world'' (ibid p.530). 

%(which is kind of Krugman's argument) 
What does he mean? Jacobs has another example that may help to illustrate the point. She attacks arguments that comparative advantage originates from the division of labour as teleological; ``one might as well say rain is beneficial to plants and that that is why it rains.'' (\citealt{jacobs_cities_1986} p.70) For Jacobs, as well as Hayek, the logic of causation is wrong. As she says, of the definition of efficiency intrinsic to increasing returns: ``to say that the economy of Uruguay... was more efficient because more specialised than the economy of Switzerland is to stand reality on its head'' (ibid p.71). Jacobs wants to build a causal argument about the role of `nests of symbiotic suppliers'; any black box approach makes this impossible. Above, however, it was argued that treating economic entities as black boxes is acceptable, depending on what argument is being made. So what is the problem? 

The reason this point has been laboured is because it transpires there are very specific ways in which Hayek and Jacobs are correct: when the problem of production in space is posed from actors' points of view, several of the key functional assumptions of constrained optimisation cannot work. Some `black boxes' can still be used, however, and the issue is not quite as clear cut as Hayek's job-lot dismissal would imply. The devil is in the detail, however.

---

%%%%%%%%%%%
% Brute force vs analytic
%%%%%%%%%%%
%As discussed, the two-good models use an `Engel function' able to either act as a Cobb Douglas if $\gamma$ is set to zero, or to distinguish the two goods as necessity and luxury (see section \ref{sec:utilitynengelslaw}) . 
%In an analytic setting, using  constraints (equations \eqref{eq:engelsmurataFopt} and \eqref{eq:engelsmurataFopt}) would be normal. 
As discussed above, for a given budget and set of prices, equations can be deduced that say precisely what the optimal quantity of goods should be. In an agent setting, it is possible to find optima through brute computational force: simply search the entire parameter space and keep whichever produces the highest utility. This is, of course, more computationally expensive than using a simple constrained optimisation equation, so it is a natural step to wonder whether those equations can be used instead of a brute force approach. The answer, surprisingly, is no. 
%Section \ref{sec:blackbox_tel} picked up on some of Hayek and Jacobs' arguments that gave a background to why this might be so. Here, some modelling reasons that support their argument are discussed.

Consider three options for an actor finding its maximum utility - brute force, or one of two constrained optima, the first constrained by a given total quantity of the two goods, the second including the price of both goods. In the brute force approach, the actor searches a range of possible mixes of $F$ and $M$ in Cobb Douglas utility function \eqref{eq:-3} to find the optimal mix:

\begin{equation}
max\; U=F^{\delta}M{}^{1-\delta},\; (0<\delta<1),\; F+M=Y
%U=F^{\delta}(M+\gamma){}^{1-\delta},\; (0<\delta<1),\; \gamma\geq0,\; F+M=Y\label{eq:engelsmurata2}
\end{equation}

As discussed earlier, for a budget of $F+M=Y$, the non-price constrained optima for the two goods is:

\begin{equation}
F=\delta Y
%F=\delta(Y+\gamma)\label{eq:engelsmurataFopt2}
\end{equation}

\begin{equation}
M=(1-\delta) Y
%M=Y-\delta(Y+\gamma)\label{eq:engelsmurataMopt2}
\end{equation}

Including price, the optima are as follows ($p_F$ and $p_M$ are the price of $F$ and $M$ respectively).

\begin{equation}
F=\frac{\delta Y}{p_F}
%M=\frac{Y(1-\delta)}{p_{M}}-\delta\gamma\label{eq:engelpriceoptM}
\end{equation}
 
\begin{equation}
M=\frac{(1-\delta) Y}{p_M}
%F=\frac{\delta Y-Y+\delta\gamma p_{M}}{p_{F}}\label{eq:engelpriceoptF}
\end{equation}

So why can't either two of these constrained optimisations be used? It is impossible to assume that goods will be available to meet a set `objective' demand. Consider what people must do (outside of the choice of backyard, which avoids all co-ordination problems.) They need to find information. Either -

\begin{itemize}
	\item How much of F or M can I get for my time?
	\item How much money can I get for my time, and how much of F and M will that buy?
\end{itemize}

Even if time or money can be used to produce an `objective demand', the model probably cannot supply it. The only option for people is to check what bundles of goods are available and test which mix is best. Compare this to equations \ref{eq:engelsmurataFopt2} and \ref{eq:engelsmurataMopt2}: both require prior knowledge of the budget, $Y$ - that is, the total quantity of goods. That `budget level' is not given - there is a wage, but not an `amount of $F$ plus $M$'. 


%%%%%%%%%%%%%%%%%%
% BLACK BOXES 
%%%%%%%%%%%%%%%%%%
\subsection{`Black boxes' in production and utility \label{sec:blackboxes}}

Er, what was I going to say here?
%%%%%%%%%%%%%%%%%%%%%%%%%%
% CUTTINZ
%%%%%%%%%%%%%%%%%%%%%%%%%%

%While Solow and Krugman are in large part correct, there are problems  with their case. Most fundamentally, there is a large gap between, say, Krugman's stated awareness of a model's ``silly assumptions'' and the uses to which those models are now being put in policy circles. 

%Pursuit of that argument, however, is beyond the scope of the thesis. The more important issue at hand is this: there are implications for agent modelling. How so? It comes back to Daly's claim that standard economic functions offer only 	 (p.262) Consider an agent model of production that has aspirations to examine the complex morphology of linkages seen in reality - Jane Jacobs' `tangled bank'. Agent models are almost always explicitly time-driven; computer cycles are synonymous with timesteps and iteration is central. In such a world, what moves between economic linkages each timestep, if not actual resources? Certainly, many exchanges can be service-related, but in a time-driven model it becomes much harder to conceive of resources in an abstract away. Instead, a more Weberian image presents itself as the most likely: resources with weights, combined to produce other resources with weights.

%Trying to think this through. It seemed important to me because I can't get utility functions to work in an agent context. But then in what circumstances am I saying its OK to use increasing returns? This might be better in that section.
%Her answer: "economic forces originating in import-replacing cities shape economic life outside of cities... when the force at issue is the power of a city market bearing upon regions that lack import-replacing cities of their own, the results are specialised regions." [70]
%She also argues with the definition of efficiency intrinsic to narrow EoS: "to say that the economy or Uruguay, 'the Switzerland of South America', was more efficient because more specialised than the economy of Switzerland is to stand reality on its head." [71]
%And, of course, it's right  - I think. I'm just mulling over the weird role of the theory of comparative advantage. It purports to describe, but it also underpins a prescription: opening up to trade capitalises on comparative advantage. But where does it come from in the first place? In krugman, it's just a function of input. But - as Jacobs rightly points out - that tells us nothing about its actual origins.

%%%%%%%%%%%%%%%%%%%%%%%%%%
% CRITIQUES N TING
%%%%%%%%%%%%%%%%%%%%%%%%%%
%Jane Jacobs' arguments present a good critique for the straightforward  'black box' model. Where the point of those models is that an in-out relationship can be defined - and only altered through some technological change - Jacobs' description of production paints a picture of economic actors stumbling upon new combinations of goods in a way very reminiscent of evolution. 

%Perhaps the most challenging productive factor to theorise is knowledge. An example from Krugman illustrates the sort of interactions that can take place. He argues that as energy became cheaper and more accessible, it could substitute for skilled labour. The result was a replacement of that labour by cheap energy sources - and an increase in the footloose mobility of workers, because pools of local skilled labour were now less necessary. 

%Structural vs input factors


